\documentclass[10pt,a4paper]{amsart}
\usepackage[margin=19mm]{geometry}
\usepackage{amsmath,amssymb,mathtools}
\usepackage[scr=rsfs,cal=euler]{mathalfa}
\usepackage{xfrac}
\usepackage[unicode,hidelinks,bookmarks=false]{hyperref}
\newcommand{\ie}{i.e.\ }
\newcommand{\cf}{cf.\ }
\newcommand{\resp}{respectively\ }
\newcommand{\Subsection}{\subsection}
\providecommand{\nobreakdash}{\nobreak\hbox{-}}
\setlength{\emergencystretch}{2em}
\multlinegap0pt
\usepackage{bm}
\usepackage{bbm}
\usepackage{dsfont}
\usepackage{xspace}
\usepackage{cancel}
\usepackage{mathtools}
\usepackage{amsthm}
\makeatletter
\@ifundefined{corollary}{\newtheorem{corollary}{Corollary}}{}
\@ifundefined{lemma}{\newtheorem{lemma}[corollary]{Lemma}}{}
\@ifundefined{theorem}{\newtheorem{theorem}[corollary]{Theorem}}{}
\makeatother
\title[Improved bounds on the zeros of the chromatic polynomial of ]{Improved bounds on the zeros of the chromatic polynomial of graphs and claw-free graphs}
\author{Ferenc Bencs, Guus Regts}
\date{Source: arXiv:2505.04366v2 (CC BY 4.0)}
\begin{document}
\maketitle

\noindent\emph{Source and licence.} Improved bounds on the zeros of the chromatic polynomial of graphs and claw-free graphs. Version arXiv:2505.04366v2 (CC BY 4.0), distributed under the Creative Commons Attribution 4.0 International licence, as recorded in the arXiv licence metadata for this exact version and shown on the abstract page https://arxiv.org/abs/2505.04366v2. Full rights notice: licenses/arxiv-2505.04366-CC-BY-4.0.txt.\par\smallskip

\noindent\textit{Editorial scope.} Counted unit: the Theorem whose source label is \texttt{thm:main general} in the article (source0.tex lines 603--621, source label \texttt{thm:main general}), delivered together with the article's own complete original proof of it (source0.tex lines 622--695, one \texttt{proof} environment of 74 lines). No mathematics is written, repaired or re-derived: statement and proof are the article's own text line for line. Required context retained uncounted: eq:deletion/contraction F (lines 270--272); lem:expand G/e (lines 560--567). Every definition, display and label of those lines is retained unchanged. Omitted from this package and not redistributed: 1--269, 273--559, 568--602, 696--1089. Nothing inside a retained excerpt is omitted or altered: every retained body is delivered exactly as the article's lines are written, so no omission receipt of any kind accompanies this package. Citations: the retained excerpts carry no citation command at all, so no bibliography entry of the article is transcribed and no reference is redistributed. Reference adaptation: none was needed and none was made; every reference inside the delivered unit resolves inside it, so no unresolved reference or citation is produced. Printed slips kept literal: no printed word, symbol, hypothesis or number of the counted statement or of its proof was altered. Layout and class: the article's own class is \texttt{amsart} with packages including amssymb, enumitem, amsrefs, tikz, blindtext, xcolor, comment, mathrsfs, geometry, a4wide, fullpage, mathtools, hyperref, graphicx; the wrapper is the lane's pinned \texttt{amsart} editorial wrapper at 10pt on A4, which restates the theorem-like environments the retained text needs and adds the article's own macro definitions that the retained text needs (none), with extra wrapper packages (bm, bbm, dsfont, xspace, cancel, mathtools). The delivered numbering is the unit's own. Assets: no retained span contains a figure environment, a graphics-inclusion command, a PDF-page inclusion, a table or a TikZ picture; no image file is copied into this package, the package declares no asset dependency and no image file is redistributed with it. Licence: version arXiv:2505.04366v2 of this article is distributed under the Creative Commons Attribution 4.0 International licence, as recorded in the arXiv licence metadata for this exact version and shown on the abstract page https://arxiv.org/abs/2505.04366v2. A full rights notice is delivered at licenses/arxiv-2505.04366-CC-BY-4.0.txt. Characters: every retained excerpt is pure ASCII, so the delivered engine, XeLaTeX with its default Latin Modern text font, reports no missing character.
\begin{equation}\label{eq:deletion/contraction F}
    F_G(x)=F_G(x)+xF_{G/e}(x).
\end{equation}

\begin{lemma}\label{lem:expand G/e}
Let $G=(V,E)$ be a graph with an ordering of the edges and let $e=\{u,v\}$ be the smallest edge. Assume that the other edges incident to $v$ are the largest edges.   
Then
\begin{equation}\label{eq:contraction recurrence}
F_{G/e}(x)=\sum_{\substack{T\in \mathcal{F}_{G\setminus e,u}\\ V(T)\cap N_{G\setminus e}(v)=\emptyset}}x^{|T|}F_{G\setminus V(T)}(x),    
\end{equation}
where $N_{G\setminus e}(v)$ denotes the collection of neighbors of $v$ in the graph $G\setminus e$.
\end{lemma}

\begin{theorem}\label{thm:main general}
Let $\Delta\geq 2$ and $g\geq 3$ be integers. 
Let $a\in (0,1)$ and $b>0$. 
Let for $i=0,\ldots,\Delta-1,$
\begin{equation}\label{eq:define a_i}
    a_i:=\frac{b\left(1-a+ t_{\Delta,g,i}(b)\right)}{(\Delta-1) \left(1+\tfrac{b}{\Delta-1}\right)^{\Delta-1}}.
\end{equation}
Suppose that
\begin{itemize}
    \item $a_i<1$ for $i=0,\ldots,\Delta-1$, and
    \item $\prod_{i=0}^{\Delta-2}(1-a_i)\geq 1-a$.
\end{itemize}
%\[
%    \prod_{i=0}^{\Delta-2} \left(1-\frac{b\left(1-a+ t_{\Delta,g,i}(b)\right)}{(\Delta-1) \left(1+\tfrac{b}{\Delta-1}\right)^{\Delta-1}}\right) \geq 1-a.
%\]
Then for any $z\in \mathbb{C}$ such that 
\[|z|\leq \frac{b(1-a)}{(\Delta-1)(1+\tfrac{b}{\Delta-1})^{\Delta-1}}\]
and any graph $G$ of maximum degree at most $\Delta$ and girth at least $g$ we have $F_G(z)\neq 0$.
\end{theorem}

\begin{proof}
%Let us define for $i=0,\ldots,\Delta-2$, 
% \[
% a_i:=\frac{b\left(1-a+ t_{\Delta,g,i}(b)\right)}{(\Delta-1) \left(1+\tfrac{b}{\Delta-1}\right)^{\Delta-1}}.
% \]
% Then the condition on $a$ and $b$ translates to $\prod_{i=0}^{\Delta-2}(1-a_i)\geq 1-a.$
Fix $z$ such that $|z|\leq \frac{b(1-a)}{(\Delta-1)(1+\tfrac{b}{\Delta-1})^{\Delta-1}}$.
We will prove the the following three statements by induction on the number of edges:
\begin{itemize}
    \item[(i)] for any edge $e$ and any vertex $u\in e$, $\left|\frac{F_G(z)}{F_{G\setminus e}(z)}-1\right|\leq a_{\deg(u)-1}$,
    \item[(ii)] $F_G(z)\neq 0$,
    \item[(iii)] for any vertex $u$ of degree at most $\Delta-1$, $\left|\frac{F_{G-u}(z)}{F_{G}(z)}\right|\leq \frac{1}{1-a}$.
\end{itemize}
Clearly, if $G$ consists of a single edge, all three statements are true, since then $\tfrac{F_G(z)}{F_{G\setminus e}}=1+z$ and so (i) follows by assumption since $a_0=\frac{b(1-a)}{(\Delta-1)(1+\tfrac{b}{\Delta-1})^{\Delta-1}}$.
Part (ii) follows since $a_0<1$ and (iii) follows since 
\[
\left |\frac{F_{G-u}(z)}{F_G(z)}\right |=\left |\frac{1}{1+z}\right |\leq \frac{1}{1-a_0}\leq \frac{1}{1-a},
\]
by assumption.

Now assume that $G$ is a graph of maximum degree at most $\Delta$ and girth at least $g$ with more than one edge.
Note that item (i) implies item (ii), since by induction $F_{G\setminus e}(z)\neq 0$ for any edge $e$ and all $a_i$ are contained in $(0,1)$.

Secondly, statements (i) and (ii) imply statement (iii).
Indeed, by (ii) $F_G(z)\neq 0$. Write $e_1,\ldots,e_d$ for the edges incident with $u$ and denote $G_d=G$ and $G_i=G_{i+1}\setminus e_{i+1}$ for $i=0,\ldots,d-1$.
Note that $G_0=G-u$ and that the degree of $u$ in $G_{i}$ is equal to $i.$
Then 
\begin{equation}\label{eq:telescope from vertex to edge}
 \frac{F_{G-u}(z)}{F_{G}(z)}=\prod_{i=1}^{d} \frac{F_{G_{i}\setminus e_i(z)}}{{F_{G_{i}}}(z)}. \end{equation}
By (i) and induction we have 
\[
\left|\frac{F_{G_i\setminus e_i}(z)}{F_{G_i}(z)}\right|\leq \frac{1}{1-a_{i-1}}.
\]
Therefore, \eqref{eq:telescope from vertex to edge} is bounded by $\prod_{i=0}^{d-1}\tfrac{1}{1-a_i}\leq \prod_{i=0}^{\Delta-2}\tfrac{1}{1-a_i}\leq \tfrac{1}{1-a}$ by the assumptions of the theorem.

We thus focus on proving item (i) for any given edge $e$ of $G$ and $u\in e$.
Let us write $e=\{u,v\}$ and choose an ordering of the edges of $G$ in which $e$ is the smallest edge and the other edges incident to $v$ are the largest edges in the ordering of the edges of $G$.
Then by the deletion-contraction recurrence~\eqref{eq:deletion/contraction F}, 
\begin{equation}\label{eq:del/con in ratio}
\frac{F_G(z)}{F_{G\setminus e}(z)}-1=z\frac{F_{G/e}(z)}{F_{G\setminus e}(z)}=z+z\frac{F_{G/e}(z)-F_{G\setminus e}(z)}{F_{G\setminus e}(z)}.
 \end{equation}
By Lemma~\ref{lem:expand G/e} the difference $F_{G\setminus e}(z)-F_{G/e}(z)$ is given by
\[
\sum_{\substack{T\in \mathcal{F}_{G\setminus e,u}\\ V(T)\cap 
N_{G\setminus e}(v)\neq \emptyset}} z^{|T|}F_{G\setminus V(T)}(z).
\]
Any tree appearing in this sum contains a path rooted at $u$ with at least $g-2\ge 1$ edges, since the girth of $G$ is at least $g$.
Therefore, using the triangle inequality, we can bound the absolute value of the left-hand side of~\eqref{eq:del/con in ratio} by
\begin{equation}\label{eq:tree gen in proof}
|z|+\sum_{\substack{T\in \mathcal{T}_{G\setminus e,u;g-2}}}|z|^{|T|+1}\left|\frac{F_{G\setminus V(T)}}{F_{G\setminus e}(z)}\right|.
\end{equation}
Fix a tree $T$ appearing in this sum and write $V(T)=\{u=u_1,\ldots,u_k\}$ such that $u_i$ has at least one neighbour in $\{u_1,\ldots,u_{i-1}\}$ for $i\geq 2$. Denote $G_0=G\setminus e$ and for $i=1,\ldots, k$ $G_i=G_{i-1}-u_i$.
Then 
\begin{equation}\label{eq:telescope}
\frac{F_{G\setminus V(T)}(z)}{F_{G\setminus e}(z)}=\prod_{i=1}^k \frac{F_{G_{i-1}-u_i}(z)}{F_{G_{i-1}}(z)}.
\end{equation}
Since the degree of $u_i$ in $G_{i-1}$ is at most $\Delta-1$, we can apply induction and (iii) to bound 
\[
\left|
\frac{F_{G\setminus V(T)}(z)}{F_{G\setminus e}(z)}\right|\leq (1-a)^{-|V(T)|}.
\]
%Using Lemma~\ref{lem:tree} applied to the graph $G\setminus e$ to bound the BCF tree-generating function, 
We can thus bound~\eqref{eq:tree gen in proof} by
\begin{align*}
&|z|+\frac{|z|}{1-a}\sum_{T\in \mathcal{T}_{G\setminus e,u;g-2}}\left(\frac{|z|}{1-a}\right)^{|T|} 
\\
\leq &|z|+\frac{|z|}{1-a}T_{G\setminus e,u;g-2}\left(\frac{b}{(\Delta-1)(1+\tfrac{b}{\Delta-1})^{\Delta-1}}\right)^{|T|}
\\
\leq &|z|+\frac{|z|}{1-a}t_{\Delta,g,\deg(u)-1}(b)= a_{\deg(u)-1},
%\label{eq:tree gen in proof step 2}
\end{align*}
as desired.
This finishes the proof.
\end{proof}

\end{document}
